% =====================================================================
%  RESUMO EM LÍNGUA PORTUGUESA (obrigatório) -- Manual, seção 2.1.8
% ---------------------------------------------------------------------
%  Regras (ABNT NBR 6028:2021):
%   - UM ÚNICO parágrafo, sem recuo, espaço simples;
%   - entre 150 e 500 palavras;
%   - frases concisas e afirmativas; terceira pessoa do singular e voz
%     ativa ("Este trabalho analisa...", e não "Neste trabalho eu analisei...");
%   - deve apresentar: OBJETIVO, MÉTODO, RESULTADOS e CONCLUSÕES;
%   - sem citações, figuras, tabelas ou equações;
%   - palavras-chave logo abaixo, separadas por ponto e vírgula (;) e
%     terminadas com ponto final (.).
%  DICA: escreva o resumo por último, quando o trabalho estiver pronto.
% =====================================================================
\setlength{\absparsep}{18pt}
\begin{resumo}
\SingleSpacing
\noindent A evasão de praticantes é um problema recorrente em academias de
lutas e artes marciais, e muitos alunos abandonam o treino antes de
perceber os seus benefícios físicos e mentais. Este trabalho apresenta o
RoundUp, um aplicativo móvel gamificado que busca aumentar o engajamento e
a retenção dos alunos e apoiar professores e academias na gestão das
turmas. O projeto foi conduzido como um processo de ideação do autor:
treze ideias de funcionalidades foram convertidas em requisitos,
construídas como hipóteses e, em paralelo, submetidas a um questionário
que serviu para confirmá-las ou refutá-las e para revelar ideias novas.
O questionário foi respondido por praticantes e ex-praticantes reais e,
para cobrir perfis ausentes, como o de professor, por 65 respostas
sintéticas, geradas por um \textit{script} a partir de 12 personas
elaboradas com apoio de inteligência artificial; essas respostas
trouxeram ideias novas e um sinal complementar, mas a decisão sobre cada
requisito se apoiou nos dados reais. O protótipo foi desenvolvido em
React Native com Firebase, com três perfis de usuário (aluno, professor e
academia), e implementa os 33 requisitos funcionais, entre eles registro
de presença com validação pelo professor, pontos de experiência, níveis,
sequência de treinos, medalhas, classificação, missões, registro de
graduação, relatórios de frequência e consentimento conforme a LGPD. A
maior parte do código foi escrita por agentes de IA, enquanto o autor
concentrou o seu trabalho no planejamento, no projeto e no controle de
qualidade dos agentes e do software em si: as regras de negócio que protegem a gamificação foram
implementadas nas regras de segurança do banco de dados e cobertas por
288 testes automatizados, e quatro requisitos foram verificados em
aparelho. Na avaliação real, seis das doze ideias avaliadas foram
confirmadas, com destaque para os recursos que tornam o progresso
visível, como a barra de progresso até a graduação e o histórico de
presença; lembretes, mensagens automáticas, missões e o próprio registro
de presença não foram confirmados, e os comentários apontaram o excesso
de notificações e de dados a inserir como riscos de abandono, já tratados
por requisitos do projeto. As respostas sintéticas e os comentários
resultaram em nove requisitos candidatos para versões futuras. Conclui-se
que a gamificação centrada no progresso visível é a direção mais
promissora para o aplicativo, que parte da evasão tem causas externas que
um aplicativo não resolve, que o desenvolvimento com agentes de IA
desloca o papel do desenvolvedor para a especificação e a verificação, e
que a próxima etapa é validar o aplicativo e os requisitos de gestão com
alunos e professores em campo.

\vspace{\onelineskip}
\noindent\textbf{Palavras-chave}: gamificação; artes marciais; engenharia de requisitos; personas; desenvolvimento assistido por IA.
\end{resumo}